% GENERATED FILE. Edit canonical lessons, then run npm run generate:books.
% canonical-source-hash: fnv1a64:b4bef7c132d665e3
% canonical-lessons: AR-C40-limadha, AR-C40-lianna, AR-C40-anna, AR-C40-hinama, AR-R40-why-that-when

\chapter{Why, That, and When}
\label{ch:ar-why-that-and-when}
\hlchaptermodality{40}

\begin{chapteropening}
\textbf{By the end of this chapter:} \emph{I can ask why with \ar{لماذا}, answer it with \ar{لأن}, report what I write or think with \ar{أنّ}, and put one event beside another in time with \ar{حينما}.}

The first three-clause Arabic sentence in this track: an acceptance, an objection and a reason, built from \ar{حسنًا}, \ar{لكن} and \ar{لأن} over the present tense of chapter 37.
\end{chapteropening}

\section[limādhā?]{\ar{لماذا} (limādhā) — why}

\label{lesson:AR-C40-limadha}



\begin{quote}
This book can ask \emph{what}, \emph{how}, \emph{how much} and — since the last chapter — \emph{yes or no}. It has never been able to ask for a reason.

\begin{itemize}
\raggedright
  \item \emph{Recall:} the particle from the lesson before, and what it moves — \emph{hal}, and nothing
\end{itemize}
\end{quote}

\subsection*{You'll want to know: \ar{لماذا}}
\begin{quote}
\textbf{\ar{لماذا؟}} — \emph{limādhā?} — \textbf{why?}
\end{quote}

It is a compound and every piece is recoverable: \textbf{\ar{لـ}}, the preposition \emph{for}, plus \textbf{\ar{ماذا}}, \emph{what} — which is itself \textbf{\ar{ما}} plus \textbf{\ar{ذا}}, \emph{this}. Arabic's \emph{why} is transparently \textbf{for-what-this}.

And the \textbf{\ar{ما}} inside it is the same particle the last chapter taught as the negator of the past tense, and the same one chapter two taught as \emph{what} in \textbf{\ar{ما} \ar{اسمك}}. Three jobs, one word, told apart by where it stands.

\begin{cousinweb}
\textbf{\ar{لماذا}} and \textbf{\ar{هل}} are the two ends of Arabic questioning. \textbf{\ar{هل}} asks for a yes or a no and adds no information; \textbf{\ar{لماذا}} asks for a whole clause back, and the next lesson supplies the word that clause begins with.

The letter \textbf{\ar{ذال}} in the middle of this word has no lesson in this track's script ladder yet, so the word is printed whole here, in the lesson that owns it, and named in romanization wherever else it appears.
\end{cousinweb}

\subsection*{Your turn}
Say these aloud:

\begin{itemize}
\raggedright
  \item \emph{limādhā?} — on its own
  \item the question words this book now owns — \emph{mā, kayfa, kam, hal, limādhā}
  \item the pieces it is built from — \emph{li- … mā … dhā}
\end{itemize}

\begin{tcolorbox}[breakable,colback=teal!4,colframe=teal!35!black,title={Before you move on}]
What are the three pieces of \textbf{\ar{لماذا}}? (\textbf{\ar{لـ}} \emph{for}, \textbf{\ar{ما}} \emph{what}, \textbf{\ar{ذا}} \emph{this}.) Which of them has the book already taught twice, in two different jobs? (\textbf{\ar{ما}} — \emph{what} in chapter two, \emph{did not} in chapter thirty-nine.) And what does \textbf{\ar{لماذا}} ask for that \textbf{\ar{هل}} does not? (\textbf{A whole clause}.)

Source: \href{https://www.unicode.org/charts/PDF/U0600.pdf}{Unicode Arabic chart}.
\end{tcolorbox}
\section[li-anna]{\ar{لأن} (li-anna) — because}

\label{lesson:AR-C40-lianna}



\begin{quote}
The lesson before this taught the question. A question nobody can answer is worse than no question at all.

\begin{itemize}
\raggedright
  \item \emph{Recall:} the item before this, and the three pieces inside it — \emph{limādhā}
\end{itemize}
\end{quote}

\subsection*{You'll want to know: \ar{لأن}}
\begin{quote}
\textbf{\ar{لا} \ar{أكتب} \ar{لأنني} \ar{لا} \ar{أسمع}} — \emph{lā aktub li-annanī lā asmaʿ} — \textbf{I am not writing because I do not hear.}
\end{quote}

Look at the first letter. \emph{limādhā} opens with \textbf{\ar{لـ}}, \emph{for}, and so does this word — with \textbf{\ar{أنّ}}, \emph{that}, behind it instead of a question word. Arabic's \emph{why} and \emph{because} are one particle apart, and both are visibly built out of pieces rather than borrowed whole.

\textbf{\ar{لأنّ}} needs a subject after it, so for \emph{I} it becomes \textbf{\ar{لأنني}}, \emph{li-annanī}. That is one ending on one word, and it is the only thing the construction asks of you.

\begin{grammarlens}[title={Grammar Lens: the claim, then the reason}]
\noindent
\begin{tabularx}{\linewidth}{@{}>{\raggedright\arraybackslash}X>{\raggedright\arraybackslash}X@{}}
\toprule
\textbf{head word} & \textbf{what follows it} \\
\midrule
\textbf{\ar{لكن}} & an objection \\
\textbf{\ar{لأن}} & a reason \\
\bottomrule
\end{tabularx}

Both stand at the head of the clause they govern, and both come \textbf{after} the claim they attach to. Put them in one sentence and you have the shape of an argument.
\end{grammarlens}

\subsection*{Your turn}
Say these aloud:

\begin{itemize}
\raggedright
  \item the question, then the answer word — \emph{limādhā?} … \emph{li-anna}
  \item \emph{lā aktub li-annanī lā asmaʿ}
  \item an objection and a reason in one breath — \emph{ḥasanan, lākin … li-anna …}
\end{itemize}

\begin{tcolorbox}[breakable,colback=teal!4,colframe=teal!35!black,title={Before you move on}]
Which letter do \emph{limādhā} and \textbf{\ar{لأن}} share, and what does it mean? (\textbf{\ar{لـ}}, \emph{for}.) What follows it in each? (\textbf{A question word} in one, \textbf{a complementiser} in the other.) And which comes first, the claim or the reason? (\textbf{The claim}.)

Source: \href{https://www.unicode.org/charts/PDF/U0600.pdf}{Unicode Arabic chart}.
\end{tcolorbox}
\section[anna]{\ar{أنّ} (anna) — that}

\label{lesson:AR-C40-anna}



\begin{quote}
\emph{fakkara} — \emph{he thought} — has been in this book since chapter thirty-two, and nobody has ever been able to say \textbf{what} they think.

\begin{itemize}
\raggedright
  \item \emph{Recall:} the reason word from the lesson before, and the piece behind its \ar{لـ} — \emph{li-anna}
\end{itemize}
\end{quote}

\subsection*{You'll want to know: \ar{أنّ}}
\begin{quote}
\textbf{\ar{أكتب} \ar{أنّ} \ar{سامي} \ar{جاري}} — \emph{aktub anna Sāmī jārī} — \textbf{I write that Sami is my neighbour.}
\end{quote}

\textbf{\ar{أنّ}} takes a whole sentence and makes it the object of a verb of writing, saying, knowing or thinking. Everything before it is the frame; everything after it is the reported thing.

You met it one lesson ago without being told: \textbf{\ar{لأنّ}} is \textbf{\ar{لـ}} plus this word. The shadda over the \textbf{\ar{ن}} is the doubling mark, and it belongs to the word — this is \textbf{anna}, with the \emph{n} doubled, not \textbf{ana}.

\begin{grammarlens}[title={Grammar Lens: what follows it}]
The clause after \textbf{\ar{أنّ}} is a \textbf{nominal} sentence — two nouns side by side, the shape chapter two taught — and \textbf{\ar{أنّ}} puts the first of them in the accusative, in the same way \textbf{\ar{ليس}} does to what follows it. Two particles, one grammatical habit, and it is worth seeing them together.

English lets \emph{that} be dropped: \emph{I think it is good}. Arabic does not.
\end{grammarlens}

\subsection*{Your turn}
Say these aloud:

\begin{itemize}
\raggedright
  \item \emph{aktub anna Sāmī jārī}
  \item the word it hides in — \emph{li-anna} … \emph{anna}
  \item with the verb from chapter thirty-two, in romanization — \emph{fakkara anna …}
\end{itemize}

\begin{tcolorbox}[breakable,colback=teal!4,colframe=teal!35!black,title={Before you move on}]
Where have you already read \textbf{\ar{أنّ}}? (\textbf{Inside \ar{لأنّ}}, in the lesson before.) What kind of clause follows it? (\textbf{A nominal sentence}, whose first noun goes into the accusative.) And can it be dropped, the way English drops \emph{that}? (\textbf{No}.)

Source: \href{https://www.unicode.org/charts/PDF/U0600.pdf}{Unicode Arabic chart}.
\end{tcolorbox}
\section[ḥīnamā]{\ar{حينما} (ḥīnamā) — when}

\label{lesson:AR-C40-hinama}



\begin{quote}
\textbf{\ar{أنّ}} hangs a sentence off a \textbf{verb}. This lesson's word hangs one off a \textbf{moment}.

\begin{itemize}
\raggedright
  \item \emph{Recall:} the complementiser from the lesson before — \emph{anna}
\end{itemize}
\end{quote}

\subsection*{You'll want to know: \ar{حينما}}
\begin{quote}
\textbf{\ar{حينما} \ar{جاء،} \ar{كتب}} — \emph{ḥīnamā jāʾa, kataba} — \textbf{when he came, he wrote.}
\end{quote}

\begin{quote}
\textbf{\ar{أكتب} \ar{حينما} \ar{أسمع}} — \emph{aktub ḥīnamā asmaʿ} — \textbf{I write when I hear.}
\end{quote}

The clause it opens can come first or second, and Arabic marks the seam with a comma when it comes first — which is the same rule English follows and one of the few places the two languages punctuate alike.

\begin{cousinweb}
\textbf{\ar{حين}} is an ordinary noun meaning \textbf{a time, a moment}, and \textbf{\ar{ما}} is the same general particle you have now met three times: as \emph{what} in chapter two, as the past-tense negator in chapter thirty-nine, and inside \emph{limādhā} in this chapter. Stuck onto \textbf{\ar{حين}} it makes \emph{at-the-moment-that}.

Arabic writes another word here, \emph{ʿindamā}, built the same way on a preposition meaning \emph{at}. It is the commoner of the two in modern prose and it needs a letter this track's script ladder has not reached, so it is named in romanization and left for a script lesson. \textbf{\ar{حينما}} costs no new letters at all.
\end{cousinweb}

\subsection*{Your turn}
Say these aloud:

\begin{itemize}
\raggedright
  \item \emph{ḥīnamā jāʾa, kataba}
  \item the same joined the other way round — \emph{kataba ḥīnamā jāʾa}
  \item \emph{aktub ḥīnamā asmaʿ}
\end{itemize}

\begin{tcolorbox}[breakable,colback=teal!4,colframe=teal!35!black,title={Before you move on}]
What are the two pieces of \textbf{\ar{حينما}}? (\textbf{\ar{حين}}, a moment, and \textbf{\ar{ما}}.) How many jobs has that \textbf{\ar{ما}} now done in this book? (\textbf{Three} — \emph{what}, \emph{did not}, and this one.) And when does Arabic put a comma at the seam? (\textbf{When the time clause comes first}, as English does.)

Source: \href{https://www.unicode.org/charts/PDF/U0600.pdf}{Unicode Arabic chart}.
\end{tcolorbox}
\section[limādhā · li-anna · anna · ḥīnamā]{Why, that, and when}

\label{lesson:AR-R40-why-that-when}



\begin{quote}
Four items landed in this chapter and three of them share a piece with another. Name all four, and say which piece each one shares.
\end{quote}

\subsection*{Your turn}
\begin{itemize}
\raggedright
  \item \emph{Say it:} the question and its answer — \emph{limādhā?} … \emph{li-anna …}
  \item \emph{Say it:} a reported thought — \emph{aktub anna Sāmī jārī}
  \item \emph{Say it:} a moment — \emph{ḥīnamā jāʾa, kataba}
  \item \emph{Say it:} the three-clause sentence — \emph{ḥasanan, lākin lā afham li-annanī lā asmaʿ}
  \item \emph{Recall:} six lessons back, the particle that makes any statement a question — \emph{hal}
  \item \emph{Recall:} eleven lessons back, the joiner written with no space after it — \emph{wa-}
\end{itemize}

\begin{tcolorbox}[breakable,colback=teal!4,colframe=teal!35!black,title={Before you move on}]
Which two of this chapter's words begin with the same \textbf{\ar{لـ}}? (\emph{\textbf{limādhā}} and \textbf{\ar{لأن}}.) Which particle sits inside \textbf{\ar{لأن}}? (\textbf{\ar{أنّ}}.) And which of the four hangs a clause off a moment rather than off a verb? (\textbf{\ar{حينما}}.)
\end{tcolorbox}
